% q3.tex — 六、问题三的模型建立与求解（2026-08-15, 素材见 PAPER_HANDOFF.md §7）

\subsection{数据准备}

问题三的候选单品为 2023 年 6 月 24 日至 30 日有销售记录的 49 种蔬菜单品。对每个
候选单品计算五个评价指标：周利润额（销售额扣除按批发价与损耗折算的成本）、周销售
量、周销售次数（交易笔数）、周打折次数与损耗率。此外，为后续优化模型准备单品级
参数：批发价取 6 月 24 日至 30 日均值，损耗率取自附件4，单品需求曲线用近 30 天
日销量对日售价的线性回归估计，样本不足 5 天的单品按"其历史销量占所属品类销量的
份额"缩放品类需求曲线（避免个别样本稀少的单品获得失真的大需求，具体修正效果在
6.5 节对比中体现）。

\subsection{基于熵权TOPSIS 的候选单品评价}

选品不能只看利润：利润额大的单品可能伴随高打折频次与高损耗，销售次数稀少的单品
则缺乏稳定的利润贡献。因此建立五项指标的评价体系，方向与熵权法求得的权重见表7。

\begin{table}[H]
\centering
\caption{候选单品评价指标与熵权}
\label{tab:topsis_ind}
\begin{tabular}{lcc}
\toprule
指标 & 指标方向 & 熵权 \\
\midrule
周利润额 & 极大型 & 0.291 \\
周销售量 & 极大型 & 0.304 \\
周销售次数 & 极大型 & 0.249 \\
周打折次数 & 极小型 & 0.119 \\
损耗率 & 极小型 & 0.037 \\
\bottomrule
\end{tabular}
\end{table}

熵权法由数据自身的离散程度确定权重。对第 $j$ 项指标，先做正向化与归一化得到
$r_{ij}$，再计算信息熵

\begin{equation}
e_j=-\frac{1}{\ln n}\sum_{i=1}^{n}p_{ij}\ln p_{ij},\qquad
p_{ij}=\frac{r_{ij}}{\sum_{i=1}^{n}r_{ij}},
\label{eq:entropy}
\end{equation}

其中 $n$ 为候选单品数。指标离散程度越大，$e_j$ 越小，信息量越大，权重
$w_j=(1-e_j)/\sum_j(1-e_j)$ 越高。五项指标中利润额、销售量与销售次数的权重
合计 0.844，说明销量规模类指标在选品中占主导。

TOPSIS 以加权归一化矩阵 $v_{ij}=w_j r_{ij}$ 为基础，取各指标最优值构成正理想解
$V^+$，最劣值构成负理想解 $V^-$，计算各单品到两者的欧氏距离：

\begin{equation}
D_i^+=\sqrt{\sum_{j=1}^{5}(v_{ij}-V_j^+)^2},\qquad
D_i^-=\sqrt{\sum_{j=1}^{5}(v_{ij}-V_j^-)^2},
\label{eq:topsis_d}
\end{equation}

相对贴近度

\begin{equation}
S_i=\frac{D_i^-}{D_i^++D_i^-}
\label{eq:topsis_s}
\end{equation}

度量单品相对最优方案的接近程度，$S_i$ 越大越值得上架。得分前五的单品为小米椒(份)
（0.744）、云南生菜(份)（0.628）、西兰花（0.585）、云南油麦菜(份)（0.500）与
芜湖青椒(1)（0.488），完整排名见支撑材料。取前 33 名单品构成优化模型的候选池，
同时保留模型在池内进一步取舍的空间。

\subsection{单品补货与定价的 0-1 规划模型}

\subsubsection{决策变量与目标函数}

对候选池内的每个单品 $i$，引入上架变量 $y_i\in\{0,1\}$，补货量 $x_i$ 与售价
$P_i$；售价在历史价格区间内离散为 $K=10$ 档候选价 $P_i^{(k)}$，用选价变量
$z_i^{(k)}\in\{0,1\}$ 表示（每单品恰选一档，$\sum_k z_i^{(k)}=y_i$）。目标为
7 月 1 日的总利润最大：

\begin{equation}
\max\ Z=\sum_{i}\left(P_i-\hat w_i\right)x_i,\qquad
\hat w_i=\bar w_i(1+\ell_i),
\label{eq:q3obj}
\end{equation}

其中 $\bar w_i$ 为近一周批发价均值。该目标未计入损耗品的打折回收，属保守口径，
实际利润不低于此估计。

\subsubsection{约束条件}

\begin{enumerate}[label=（\arabic*）]
\item 上架数量约束：$27\le\sum_i y_i\le33$；
\item 最小陈列量约束：$x_i\ge 2.5\,y_i$，上架单品至少陈列 2.5 kg；
\item 需求上限约束：$x_i\le D_i(P_i)$，$D_i$ 为单品需求曲线在所选价格处的取值，
      保证补货不超出该价格下的预期需求；
\item 价格区间约束：$P_i\in[P_i^{\min},P_i^{\max}]$，区间取 2023 年 6 月的历史
      售价范围，避免使用过时的极端价格档位；
\item 空间约束：$\sum_i x_i\le V$，$V=565.5$ kg 与问题二同口径；
\item 容量护栏：$x_i\le 1.5\,B_i^{\max}$，$B_i^{\max}$ 为单品全期最大日销量，
      防止个别单品的需求估计被异常放大；
\item 品类需求满足率下限：对每个品类 $c$，
      $\sum_{i\in c}x_i\ge 0.5\sum_{i\in c}\bar D_i$，
      其中 $\bar D_i$ 为单品 $i$ 在其 2023 年 6 月实际售价下的需求，
      该约束落实题面"尽量满足市场对各品类蔬菜商品需求"的要求。
\end{enumerate}

\subsection{求解与结果}

模型为混合整数规划，用 MATLAB 求解，耗时秒级。结果为：上架 30 个单品，7 月 1 日
总利润 1,052.95 元，空间占用 341.6 kg（容量 565.5 kg 未用满）。表8 列出单位利润
最高的 8 个单品，图13 为前 15 名单品的利润分布。

\begin{table}[H]
\centering
\caption{7 月 1 日部分单品补货量与定价（按单位利润排序的前 8 名）}
\label{tab:q3items}
\small
\begin{tabular}{llccc}
\toprule
单品名称 & 类别 & 补货量(kg) & 定价(元/kg) & 单位利润(元/kg) \\
\midrule
长线茄 & 茄类 & 8.46 & 16.00 & 8.54 \\
西峡花菇(1) & 食用菌 & 5.09 & 24.00 & 6.72 \\
西兰花 & 花菜类 & 13.83 & 15.20 & 6.65 \\
小米椒(份) & 辣椒类 & 18.12 & 8.90 & 6.56 \\
木耳菜 & 花叶类 & 4.39 & 10.00 & 6.54 \\
竹叶菜 & 花叶类 & 13.49 & 8.00 & 5.37 \\
红薯尖 & 花叶类 & 6.61 & 8.80 & 5.33 \\
净藕(1) & 水生根茎类 & 7.30 & 16.00 & 4.66 \\
\bottomrule
\end{tabular}
\end{table}

\begin{figure}[H]
\centering
\includegraphics[width=0.68\textwidth]{figures/q3_item_profit.png}
\caption{7 月 1 日选品单品利润分布（前 15 名）}
\label{fig:q3_items}
\end{figure}

各品类需求满足率见表9。六品类满足率全部不低于 50\%，其中花叶类与花菜类超过
100\%，即补货量高于其在历史实际售价下的需求估计，这是因为空间容量富余时，模型
按所选档位价格对应的需求补货，可在部分品类上适度超配。

\begin{table}[H]
\centering
\caption{7 月 1 日各品类需求满足率}
\label{tab:q3sat}
\begin{tabular}{lcccccc}
\toprule
品类 & 花叶类 & 花菜类 & 水生根茎类 & 茄类 & 辣椒类 & 食用菌 \\
\midrule
满足率 & 100.0\% & 106.6\% & 76.2\% & 103.0\% & 80.4\% & 114.8\% \\
\bottomrule
\end{tabular}
\end{table}

为考察下限取值对结果的影响，将品类满足率下限从 0 逐档提高到 50\% 重解模型，各档
结果完全相同（利润 1,052.95 元，满足率不变），说明在修正后的单品需求模型下，利润
最优解本身就兼顾了品类均衡，50\% 下限不产生额外代价，模型对"尽量满足需求"的要求
具有内在兼容性。

\subsection{与问题二的结果对比}

将问题二的品类方案与问题三的单品方案在 7 月 1 日进行对比。品类方案（方案 A 口径，
含打折回收）当日利润约 1,237 元；单品方案利润 1,052.95 元，低约 15\%。差距来自
两处：单品方案的目标采用不含打折回收的保守口径；单品需求曲线经份额缩放修正后
比品类曲线保守，且单品层面多了最小陈列量与需求满足率约束。相应地，单品方案获得
了品类方案不具备的品类均衡保障：各品类需求满足率全部不低于 50\%，且选品经过五项
指标的权衡，兼顾了销售稳定性。这一"以约 15\% 的利润换取品类均衡与选品稳健性"
的权衡，正是题目要求将补货决策细化到单品层面的实际含义。
